% archon:covers Smale5/AnalyticEnergy.lean Smale5/Diagram3NormCertificate.lean Smale5/Edges.lean Smale5/Configuration.lean Smale5/Products.lean Smale5/Balance.lean Smale5/ArchonVirial.lean Smale5/QuadraticRecovery.lean Smale5/FiniteFiber.lean Smale5/FourProductFiber.lean Smale5/PairInertia.lean Smale5/TangentEnergy.lean Smale5/Diagram7Inertia.lean
\chapter{Checked elementary results}

\begin{definition}[Normalized configuration]
\label{bp:configuration}
\lean{Smale5.Configuration, Smale5.IsNormalized}
\leanok
A configuration has five complex coordinates of each kind, ten inverse
distances, the two force balance systems, the ten inverse-distance equations,
and the gauge equating the first coordinate differences.
The definitions preserve the manuscript's twenty complex coordinates.
\end{definition}

\begin{lemma}[Four-product identities]
\label{bp:products}
\lean{Smale5.product_fourProducts, Smale5.complementary_fourProducts, Smale5.eighthPower_recovery}
\leanok
For the ten squared distances $s_e$, let $\Delta=\prod_e s_e$ and let $p_e$
be $s_e$ times the three disjoint-edge distances. Then
$\prod_e p_e=\Delta^4$, the complementary product of the $p_e$ equals
$s_e^2\Delta$, and $s_e^8\prod_f p_f$ is the fourth power of that
complementary product. These are exact commutative polynomial identities.
\end{lemma}

\begin{lemma}[Centre of mass]
\label{bp:center}
\uses{bp:configuration}
\lean{Smale5.normalized_centerOfMass_z, Smale5.normalized_centerOfMass_w}
\leanok
The normalized balance equations imply both mass-weighted centre identities.
The proof pairs antisymmetric terms in the finite double sums.
\end{lemma}

\begin{lemma}[Virial identity]
\label{bp:virial}
\uses{bp:configuration}
\lean{Smale5.virial_pair_sum, Smale5.normalized_potential_eq_inertia}
\leanok
For every normalized configuration, the potential equals the inertia:
$U=I$. No positivity assumption is needed. Multiply each first balance equation
by its weighted second coordinate, pair the two orientations of every edge,
and apply the inverse-distance equation. This proof was authored and checked
in the ongoing local task; the external Archon proving session never ran.
\end{lemma}

\begin{lemma}[Weighted pair-distance inertia]
\label{bp:pair-inertia}
\uses{bp:configuration, bp:center}
\lean{Smale5.weighted_pair_distance_identity, Smale5.normalized_mass_mul_inertia}
\leanok
For any configuration, total mass times inertia equals the weighted sum of
squared pair distances plus the product of the two weighted coordinate sums.
For normalized configurations the centre identities remove that product,
giving $MI=\sum_{i<j}m_i m_j s_{ij}$. This is an exact finite polynomial
identity; it supplies no asymptotic estimate by itself.
\end{lemma}

\begin{lemma}[Algebraic tangent energy]
\label{bp:tangent-energy}
\uses{bp:configuration}
\lean{Smale5.normalized_tangent_energy_critical, Smale5.normalized_tangent_virial, Smale5.normalized_tangent_energy_zero}
\leanok
At a normalized configuration, the formal directional expressions satisfy
$2dU+dI=0$ whenever the direction satisfies the linearized distance equations.
If it also satisfies the linearized first force equations, then $dU=dI$,
hence $dU=dI=0$. These are finite algebraic statements about the explicitly
defined directional expressions. Identification with analytic derivatives
and constancy on an irreducible complex curve remain separate obligations.
\end{lemma}

\begin{lemma}[Quadratic reconstruction]
\label{bp:quadratic}
\lean{Smale5.quadratic_offset_recovery, Smale5.finite_quadratic_roots, Smale5.finite_offset_families}
\leanok
For a nonzero fixed gauge difference and fixed squared distances from the
first two bodies, each coordinate offset satisfies a monic quadratic.
There are finitely many simultaneous offset choices.
\end{lemma}

\begin{theorem}[Finite fixed-distance fibre]
\label{bp:distance-fibre}
\uses{bp:configuration, bp:center, bp:quadratic}
\lean{Smale5.finite_normalized_fixed_distances}
\leanok
For positive masses and fixed nonzero squared distances, the original
normalized configuration fibre is finite. The centre identities reconstruct
the translation coordinates, and each inverse distance has finitely many
square-root choices. This replaces the manuscript's use of Witt's theorem.
\end{theorem}

\begin{theorem}[Finite four-product fibre]
\label{bp:fourproduct-fibre}
\uses{bp:products, bp:distance-fibre}
\lean{Smale5.finite_normalized_fixed_fourProducts, Smale5.infinite_normalized_has_nonconstant_fourProduct}
\leanok
For positive masses, each normalized fibre with all four-products fixed is
finite. Hence every infinite set of normalized configurations has a
four-product taking at least two distinct values. This set-theoretic statement
does not itself construct a complex algebraic curve or a meromorphic pole.
\end{theorem}


\begin{lemma}[Diagram-7 inertia decomposition]
\label{bp:d7-inertia}
\uses{bp:configuration}
\lean{Smale5.Diagram7.inertia_identity, Smale5.Diagram7.normalized_inertia_identity}
\leanok
The manuscript's diagram-7 inertia decomposition is an exact polynomial
identity: inertia equals the two wing-coordinate contributions, a crossed-force
torque about the fifth body, and explicit force-residual correction terms.
For normalized solutions the residual terms vanish. This proves the algebraic
identity only; it does not prove boundedness of chart coordinates, estimates
for analytic germs, or local asymptotic vanishing of inertia.
\end{lemma}

\begin{theorem}[Energy on actual differentiable paths]
\label{bp:analytic-energy}
\uses{bp:virial, bp:tangent-energy}
\lean{Smale5.normalized_path_energy_hasDerivAt_zero, Smale5.normalized_path_energy_constant_Icc}
\leanok
Actual coordinatewise differentiation of a normalized path gives zero derivatives
of potential and inertia. They are constant on a closed real interval when
coordinates are continuous on it and differentiable in its interior. This code
was generated with Archon/Codex and then repaired and kernel checked locally.
\end{theorem}

\begin{theorem}[Public P66 certificate, both families]
\label{bp:p66-certificate}
\lean{Smale5.Diagram3Norm.p66_factorization, Smale5.Diagram3Norm.cleared_nine_root_norm, Smale5.Diagram3Norm.rational_nine_root_norm, Smale5.Diagram3Norm.positive_family_root_form_ne_zero, Smale5.Diagram3Norm.p66_negative_decomposition, Smale5.Diagram3Norm.p66_negative_pos, Smale5.Diagram3Norm.negative_family_root_form_ne_zero}
\leanok
The 67 published integer coefficients are checked as data: their exact norm
factorization holds and both family obstructions are nonzero at every
pair of complex cubic-root labels for positive unequal parameters. For the
negative argument, a kernel-checked positive-coefficient decomposition proves
positivity. The preceding fifth-order jet identity is still a separate obligation.
\end{theorem}
